\subsection{Overview}

%the normal workflow
Merton Council currently handles parking suspensions largely through manual processes. Under Merton’s published suspension rules (effective April 2026), applicants must apply at least ten full working days in advance and pay upfront. Applications are submitted via an online form on the council's website (cite!); steps include entering personal details, specifying location and date, then pre-emptively submitting payment. The details are logged by the back-office in a SharePoint Excel spreadsheet for audit. A street sign and waiver is filled in from a template, with the waiver emailed back to the applicant, and the sign printed out for collection. Field engineers collect printed job sheets at 10:00 AM daily, to deploy the physical suspension signs at least seven days in advance of suspension. The applicant can display the parking waiver on their vehicle to prove permission. If unauthorised vehicles block the parking bay, the office can be called to dispatch an enforcement officer who issues a Penalty Charge Notice (PCN); however, the vehicle cannot be towed.

%the problems
Several concerns were raised about the current workflow, and how it creates substantial friction for both staff and applicants. Because the online form uses a free-text box without up-to-date bay data, applicants often submit vague or non-existent addresses, misjudge how many bays they need, or book over existing suspensions and utility works. Staff must manually infer door ranges,  cross-reference Google Maps or Controlled Parking Zone (CPZ) maps, and negotiate date changes or refunds when clashes occur. These maps lack accurate bay pinpointing and are not reliably updated by Traffic and Highways when new disabled, electric, or cycle bays are installed. Parking is a financially significant statuatory function for local councils including Merton, which recorded £25.32 m in parking income against £10.77 m in parking expenditure in 2025–26, yet streetspace data remains fragmented. The manual processes within tight deadlines force staff and applicants into uncertainty and reactive problem solving. 


%solutions and other councils
Other authorities and national initiatives demonstrate possibilities on how to reduce manual work. Councils such as Camden and Westminster feature a public interactive map where users select specific bays and see live suspensions or clashes before applying, and amendments and cancellations can be done at self-service The Department for Transport’s Digital Traffic Regulation Order (D-TRO) service also shows how national policy is pushing toward standardised machine-readable data to cut processing and enforcement friction. (cite!). In the context of Merton, where parking enforcement operates alongside platforms such as Taranto, opportunities include introducing an interactive map to block out live clashes and help users select bays upfront, embedding automated validation to reduce manual checks, integrating resident notifications for major suspensions, and enabling field engineers to record pre-parked vehicles digitally. (may need to trim this bit)

\subsection{Interview/direct engagement}

TODO in a meeting

\subsection{Rich picture}

TO INSERT